\scanpage{43}\mspage{42}

Pour At 1, At 2 c'est clair, pour At 3 il faut \deleted{de plus} imposer, en plus de

\uline{Sousat 1} $\mathcal F'$ est une partie fermée \inserted{non vide} de $\mathcal F$,

la condition

\uline{Sousat} 2 Si $F,G\in\mathcal F'$ et si pour \deleted{\unsure{tout}} $X\in\widetilde F$, $Y\in\widetilde G$, $X\vee Y$ existe dans $\mathcal F'$ (i.e. \deleted{$X\mathrel{\substack{\scriptscriptstyle\times\\[-.7ex]\scriptscriptstyle\times}}Y$\textsuperscript{\unsure{$*$|$\times$}} \textsuperscript{$*$}dans $\mathcal F$, alors $X\vee Y$ existe \unsure{donc|dans}, $X\vee Y\in\mathcal F'$ et $X\mathrel{\substack{\scriptscriptstyle\times\\[-.7ex]\scriptscriptstyle\times}}Y\unsure{\notin|\in}\mathcal F'$, alors $X\mathrel{\substack{\scriptscriptstyle\times\\[-.7ex]\scriptscriptstyle\times}}Y$ dans $\mathcal F$ et $X\vee Y\in\mathcal F'$} $X,Y$ majorés \unsure{ds|dans} $\mathcal F'$) alors $F\vee\unsure{C|G}\in\mathcal F'$ (NB on sait que $F\vee C$ existe dans $\mathcal F$).
% slanted
\margin{\uline{NB} C'est une condition \unsure{intrinsèque} à $(\mathcal F',\leq)$, i.e. la condition At 3 pour $\mathcal F'$.}

Les conditions At 4, At 5, At 6, At 7 \inserted{(pages 21, 27)} passent trivialement \inserted{une à une} de $\mathcal F$ à $\mathcal F'$. Pour At 7, je note seulement en passant
\[
\begin{aligned}
&\text{Si }X,Y\in M',\ \text{alors }X\mathrel{\overset{\circ}{\ll}}Y\text{ dans }\mathcal F\text{ équivaut à}\\
&\hspace{12em}X\mathrel{\overset{\circ}{\ll}}Y\text{ dans }\mathcal F'.
\end{aligned}\tag*{\makebox[0pt][r]{\makebox[\displaywidth][l]{(\unsure{$*$|M})}}}
\]

\deleted{Supposons maintenant $\mathcal F$ interprété comme un ensemble de figures (en fait, un sous-atelier) de $\mathrm{Fig}(M)$.}

\uline{Définition} Soit $\mathcal F$ un atelier au sens At 1 -- At 7. On appelle \uline{sous-atelier} de $\mathcal F$, une partie \deleted{\unsure{fermée}}\inserted{$\mathcal F'$} de $\mathcal F$ \deleted{\unsure{sat}}, satisfaisant \deleted{\unsure{à}}\unsure{.} Sousat 1 et Sousat 2 ci-dessus, i.e. telle que $1^\circ)$ $\mathcal F'$ soit une partie fermée non vide de $(\mathcal F,\leq)$ \deleted{et} (ce qui n'est autre que Sousat 1) et $2^\circ)$ $\mathcal F'$, munie \unsure{de de|de} $\leq$, $\ll$, est un atelier (ce qui équivaut à Sous-at 2), On dit que
